\documentclass[a4paper,12pt]{article}

\usepackage[brazil]{babel}
\usepackage[utf8]{inputenc}
\usepackage[T1]{fontenc}

\usepackage{geometry}
\usepackage{setspace}
\usepackage{titlesec}
\usepackage{hyperref}
\usepackage{graphicx}
\usepackage{caption}
\usepackage{subcaption}
\usepackage{fancyhdr}
\usepackage{xcolor}
\usepackage{amsmath}
\usepackage{amssymb}

%%%%%%%%%%%%%%%%%%%%%%%%%%%%%%%%%%%%%%%%%%%%%%%%%%
% These are some new commands that may be useful 
% for paper writing in general. If other new commands
% are needed for your specific paper, please feel 
% free to add here. 
%
% The currently available commands are organized in: 
% 1) Systems
% 2) Quantities
% 3) Energies and units
% 4) particle species
% 5) Colors package
% 6) hyperlink
%%%%%%%%%%%%%%%%%%%%%%%%%%%%%%%%%%%%%%%%%%%%%%%%%%

\usepackage{amsmath}
\usepackage{amssymb}
\usepackage{upgreek}
\usepackage{multirow}
\usepackage{setspace}% http://ctan.org/pkg/setspace
\usepackage{fancyhdr}
\usepackage{datetime}

% 1) SYSTEMS
\newcommand{\btc}               {\textbf{BTC}}
\newcommand{\btcspace}          {\textbf{BTC} }
\newcommand{\pow}               {\textbf{PoW}}

% 4) definition to references, biblatex and hyperlink
\usepackage[backend=bibtex, 
style=nature,  %style reference.
sorting=none,
firstinits=true %first name abbreviate
]{biblatex}

\usepackage{hyperref}
\hypersetup{
    colorlinks=true, %set "true" if you want colored links
    linktoc=all,     %set to "all" if you want both sections and subsections linked
    linkcolor=blue,  %choose some color if you want links to stand out
    citecolor= blue, % color of \cite{} in the text.
    urlcolor  = blue, % color of the link for the paper in references.
}

% 5) Tikz and figures
\usepackage{epsfig}
\usepackage{lmodern}
\usepackage{mathtools}
\usepackage[utf8]{luainputenc}
\usepackage{xspace}
\usepackage{tikz}
\usepackage{pgfplots}
\pgfplotsset{compat=newest}

\usetikzlibrary{positioning}
\usepackage{subcaption}

% 6) colors:
\usepackage{xcolor}
\definecolor{ao(english)}{rgb}{0.0, 0.5, 0.0} % dark green

% 7) Add lines numbers
%\usepackage{lineno}

% add pdf file to thesis:
\usepackage{pdfpages}

\usepackage[most]{tcolorbox}

%textmarker style from colorbox doc
\tcbset{textmarker/.style={%
        enhanced,
        parbox=false,boxrule=0mm,boxsep=0mm,arc=0mm,
        outer arc=0mm,left=6mm,right=3mm,top=7pt,bottom=7pt,
        toptitle=1mm,bottomtitle=1mm,oversize}}


% define new colorboxes
\newtcolorbox{hintBox}{textmarker,
    borderline west={6pt}{0pt}{yellow},
    colback=yellow!10!white}
\newtcolorbox{importantBox}{textmarker,
    borderline west={6pt}{0pt}{red},
    colback=red!10!white}
\newtcolorbox{noteBox}{textmarker,
    borderline west={6pt}{0pt}{green},
    colback=green!10!white}

% define commands for easy access
\newcommand{\note}[1]{\begin{noteBox} \textbf{} #1 \end{noteBox}}
\newcommand{\warning}[1]{\begin{hintBox} \textbf{Warning:} #1 \end{hintBox}}
\newcommand{\important}[1]{\begin{importantBox} \textbf{} #1 \end{importantBox}}

% --------------------------------------------------
% CORES E CAIXAS
% --------------------------------------------------

\definecolor{azul}{RGB}{30,90,150}
\definecolor{cinza}{RGB}{245,245,245}
\definecolor{verde}{RGB}{230,245,235}
\definecolor{amarelo}{RGB}{255,248,220}

\newtcolorbox{ideia}{
    colback=azul!5,
    colframe=azul,
    title=\textbf{IDEIA-CHAVE},
    fonttitle=\bfseries
}

\newtcolorbox{exemplo}{
    colback=verde,
    colframe=green!50!black,
    title=\textbf{EXEMPLO},
    fonttitle=\bfseries
}

\newtcolorbox{atencao}{
    colback=amarelo,
    colframe=orange!80!black,
    title=\textbf{ATENÇÃO},
    fonttitle=\bfseries
}

% --------------------------------------------------
% CABEÇALHO
% --------------------------------------------------

\pagestyle{fancy}

\fancyhf{}

\renewcommand{\headrulewidth}{0pt}

\fancyhead[R]{\thepage}

% --------------------------------------------------
% HYPERREF
% --------------------------------------------------

\hypersetup{
    colorlinks=true,
    linkcolor=blue,
    citecolor=blue,
    urlcolor=blue
}

% --------------------------------------------------
% BIBLIOGRAFIA
% --------------------------------------------------

\addbibresource{bibliography.bib}

\newcommand{\printingbibliography}{%
    \pagestyle{myheadings}
    \markright{}
    \sloppy
    \printbibliography[
        heading=bibintoc,
        title=Referências
    ]
    \fussy
}

% --------------------------------------------------
% CONFIGURAÇÕES
% --------------------------------------------------

\renewcommand{\baselinestretch}{1.5}

\geometry{
    a4paper,
    top=30mm,
    bottom=20mm,
    left=30mm,
    right=20mm
}

\titleformat*{\section}{\bfseries\large}
\titleformat*{\subsection}{\bfseries\normalsize}

\title{
    \textbf{\large
    PROVÃO PAULISTA SERIADO -- 1ª, 2ª E 3ª SÉRIES DO ENSINO MÉDIO}
}

\author{
    André V. Silva \\
    \href{https://andrevsilva.com/}{andrevsilva.com}
}

\date{\today}

% ==================================================
% DOCUMENTO
% ==================================================

\begin{document}

\maketitle

% ==================================================
% CINEMÁTICA
% ==================================================

\tableofcontents

\section{Cinemática}

A Cinemática é a parte da Física que descreve os movimentos dos corpos,
analisando grandezas como posição, deslocamento, velocidade e aceleração,
sem se preocupar inicialmente com as causas do movimento.

% --------------------------------------------------
\subsection{Representação gráfica do movimento}

Uma maneira muito importante de estudar o movimento de um corpo é utilizar
gráficos que relacionam as grandezas físicas com o tempo.

Os principais gráficos utilizados na Cinemática são:

\begin{itemize}
    \item posição (ou espaço) em função do tempo: $s \times t$;
    \item velocidade escalar em função do tempo: $v \times t$;
    \item aceleração escalar em função do tempo: $a \times t$.
\end{itemize}

\begin{ideia}

Os gráficos permitem visualizar como o movimento de um corpo acontece ao
longo do tempo.

A posição informa \textbf{onde o corpo está}; a velocidade informa
\textbf{como a posição está mudando}; e a aceleração informa
\textbf{como a velocidade está mudando}.

Podemos representar essa relação da seguinte maneira:

\[
\boxed{
\text{posição}
\longrightarrow
\text{velocidade}
\longrightarrow
\text{aceleração}
}
\]

\end{ideia}

% --------------------------------------------------
\subsection{Gráfico da posição em função do tempo}

O gráfico da posição em função do tempo representa a posição do corpo para
cada instante considerado.

No eixo horizontal representamos o tempo $t$ e no eixo vertical representamos
a posição $s$.

\[
s = f(t)
\]

A inclinação da curva fornece informações sobre a velocidade do corpo.

\begin{ideia}

No gráfico $s \times t$:

\begin{itemize}
    \item reta horizontal $\rightarrow$ posição constante;
    \item reta inclinada para cima $\rightarrow$ posição aumentando;
    \item reta inclinada para baixo $\rightarrow$ posição diminuindo;
    \item maior inclinação $\rightarrow$ maior módulo da velocidade.
\end{itemize}

A velocidade instantânea está relacionada à inclinação da curva:

\[
v = \frac{ds}{dt}
\]

\end{ideia}

% --------------------------------------------------
\subsection{Exemplo: posição em função do tempo}

Considere um corpo cuja posição varia segundo:

\[
s(t)=10+5t
\]

onde $s$ é medida em metros e $t$ em segundos.

Comparando essa expressão com a equação do movimento uniforme,

\[
s=s_0+vt,
\]

identificamos:

\[
s_0=10\,\text{m}
\]

e

\[
v=5\,\text{m/s}.
\]

Portanto, o corpo possui velocidade constante de $5\,\text{m/s}$.

\begin{exemplo}

Para $t=0$:

\[
s(0)=10\,\text{m}.
\]

Para $t=2\,\text{s}$:

\[
s(2)=10+5(2)=20\,\text{m}.
\]

Para $t=4\,\text{s}$:

\[
s(4)=10+5(4)=30\,\text{m}.
\]

Assim, a posição aumenta linearmente com o tempo.

\end{exemplo}

% --------------------------------------------------
\subsection{Gráfico da velocidade em função do tempo}

O gráfico da velocidade em função do tempo representa como a velocidade
do corpo varia durante o movimento.

No eixo horizontal representamos o tempo $t$ e no eixo vertical a velocidade
$v$.

\[
v=f(t)
\]

\begin{ideia}

No gráfico $v \times t$:

\begin{itemize}
    \item linha horizontal $\rightarrow$ velocidade constante;
    \item velocidade aumentando $\rightarrow$ movimento acelerado;
    \item velocidade diminuindo $\rightarrow$ movimento retardado;
    \item inclinação da curva $\rightarrow$ aceleração.
\end{itemize}

A aceleração instantânea é dada por:

\[
a=\frac{dv}{dt}.
\]

Além disso, a área sob o gráfico $v\times t$ representa o deslocamento
escalar:

\[
\boxed{\Delta s=\text{área sob o gráfico }v\times t}
\]

\end{ideia}

% --------------------------------------------------
\subsection{Gráfico da aceleração em função do tempo}

O gráfico da aceleração em função do tempo mostra como a aceleração varia
durante o movimento.

No eixo horizontal representamos o tempo $t$ e no eixo vertical a aceleração
$a$.

\[
a=f(t)
\]

\begin{ideia}

Quando a aceleração é constante, o gráfico $a\times t$ é uma reta horizontal.

Além disso, a área sob o gráfico da aceleração em função do tempo representa
a variação da velocidade:

\[
\boxed{
\Delta v=\text{área sob o gráfico }a\times t
}
\]

Portanto:

\[
v_f=v_i+\Delta v.
\]

\end{ideia}

% --------------------------------------------------
\subsection{Relação entre os três gráficos}

Os três gráficos estão diretamente relacionados.

\[
\boxed{
s(t)
\overset{\text{derivada}}{\longrightarrow}
v(t)
\overset{\text{derivada}}{\longrightarrow}
a(t)
}
\]

Também podemos fazer o caminho inverso utilizando áreas:

\[
\boxed{
a(t)
\overset{\text{área}}{\longrightarrow}
v(t)
\overset{\text{área}}{\longrightarrow}
s(t)
}
\]

Assim:

\begin{center}
\begin{tabular}{c|c}
\textbf{Gráfico} & \textbf{Informação obtida} \\ \hline
$s\times t$ & A inclinação fornece a velocidade \\[2mm]
$v\times t$ & A inclinação fornece a aceleração \\[2mm]
$v\times t$ & A área fornece o deslocamento \\[2mm]
$a\times t$ & A área fornece a variação da velocidade
\end{tabular}
\end{center}

% --------------------------------------------------
\subsection{Movimento uniforme}

No movimento uniforme, a velocidade permanece constante.

Assim:

\[
v=\text{constante}
\]

e

\[
a=0.
\]

A função horária da posição é:

\[
s=s_0+vt.
\]

Portanto:

\begin{itemize}
    \item o gráfico $s\times t$ é uma reta;
    \item o gráfico $v\times t$ é uma reta horizontal;
    \item o gráfico $a\times t$ coincide com o eixo do tempo.
\end{itemize}

% --------------------------------------------------
\subsection{Movimento uniformemente variado}

No movimento uniformemente variado, a aceleração permanece constante:

\[
a=\text{constante}.
\]

A velocidade varia linearmente com o tempo:

\[
v=v_0+at.
\]

A posição é dada por:

\[
s=s_0+v_0t+\frac{1}{2}at^2.
\]

Consequentemente:

\begin{itemize}
    \item o gráfico $s\times t$ é uma parábola;
    \item o gráfico $v\times t$ é uma reta inclinada;
    \item o gráfico $a\times t$ é uma reta horizontal.
\end{itemize}

% --------------------------------------------------
\subsection{Atenção à interpretação dos gráficos}

\begin{atencao}

Um gráfico de posição, velocidade ou aceleração \textbf{não representa
necessariamente a trajetória do corpo}.

Por exemplo, um gráfico $s\times t$ mostra como a posição varia com o tempo.
Ele não mostra o caminho geométrico percorrido pelo corpo no espaço.

Da mesma forma, um gráfico $v\times t$ não representa a trajetória do objeto:
ele representa a variação da velocidade ao longo do tempo.

\end{atencao}

% ==================================================
% SÍNTESE
% ==================================================

\section{Síntese}

A interpretação dos gráficos da Cinemática pode ser resumida da seguinte
forma:

\[
\boxed{
\begin{array}{c}
s(t) \\
\downarrow\ \text{inclinação}\\
v(t) \\
\downarrow\ \text{inclinação}\\
a(t)
\end{array}
}
\]

E, no sentido inverso:

\[
\boxed{
a(t)
\xrightarrow{\text{área}}
v(t)
\xrightarrow{\text{área}}
s(t)
}
\]

\begin{ideia}

A ideia fundamental é:

\[
\boxed{
\text{posição}
\rightarrow
\text{velocidade}
\rightarrow
\text{aceleração}
}
\]

A posição descreve onde o corpo está.

A velocidade descreve como sua posição muda.

A aceleração descreve como sua velocidade muda.

\end{ideia}

% ==================================================
% BIBLIOGRAFIA
% ==================================================

\printingbibliography

\end{document}